\chapter{Methodology}

\section{Pipeline Overview}
The system is a staged pipeline (Figure~\ref{fig:pipeline}). Contract text is windowed; a presence stage decides which of the 41 categories the contract contains; a span stage extracts the exact clause text for each detected category; a risk layer gives each clause a High/Medium/Low level; and a summarization stage rewrites the clause in plain English. The final output is a risk-prioritized report. Figure~\ref{fig:pipeline} names the stages and the models behind them; Figure~\ref{fig:dataflow} at the end of this chapter repeats the same chain with the concrete data object that crosses each boundary, which is the version to consult when reading the implementation.

\begin{figure}[ht]
\centering
\begin{tikzpicture}[
  node distance=4mm and 6mm,
  stage/.style={draw, rounded corners=2pt, minimum height=9mm, minimum width=30mm,
                align=center, font=\small, fill=blue!6},
  added/.style={stage, fill=orange!15},
  arr/.style={-{Stealth[length=2.5mm]}, thick}
]
\node[stage] (a) {Contract\\text};
\node[stage, right=of a] (b) {Windowing \&\\preprocessing};
\node[stage, right=of b] (c) {41-clause presence\\(TF--IDF + DistilBERT)};
\node[added, below=of c] (d) {Risk classification\\High / Med / Low};
\node[stage, left=of d] (e) {Span extraction\\(DistilBERT-QA)};
\node[added, left=of e] (f) {Plain-English\\summary (FLAN-T5)};
\node[added, below=of e] (g) {Risk-prioritized report};
\draw[arr] (a) -- (b);
\draw[arr] (b) -- (c);
\draw[arr] (c) -- (d);
\draw[arr] (d) -- (e);
\draw[arr] (e) -- (f);
\draw[arr] (f.south) |- (g.west);
\draw[arr] (e) -- (g);
\end{tikzpicture}
\caption{The end-to-end pipeline. Shaded (orange) stages --- risk classification, plain-language summarization, and the prioritized report --- are the layers this work adds beyond CUAD's original clause-extraction task.}
\label{fig:pipeline}
\end{figure}

One principle runs through everything in this chapter: \textbf{build a strong, cheap baseline first, and measure feasibility ceilings before spending training time}. Both habits changed the design in concrete ways, as the next sections show.

\section{Stage 2A: Presence Classification}

\subsection{A strong baseline first: full-document TF--IDF}
Classical bag-of-words models have no input-length limit, so a TF--IDF representation gets to read the \emph{entire} contract --- an advantage no windowed transformer has. We fit a TF--IDF vectorizer (20{,}000 features, word uni- and bi-grams, sublinear TF, $\text{min\_df}=2$) on the 408 training contracts and train one class-balanced logistic regression per category. Legal drafting is extremely formulaic (``\emph{This Agreement shall be governed by the laws of the State of \ldots}''), so distinctive vocabulary alone carries a lot of signal. This baseline sets the bar every neural model has to beat.

\subsection{Measuring the retrieval ceiling --- and dropping retrieval}
\label{sec:ceiling}
The obvious way to hand a long contract to a transformer is retrieval: slide 2{,}000-character windows with a 1{,}500-character stride over the text (the windowing scheme itself is drawn in Figure~\ref{fig:windows}), score each window by TF--IDF cosine similarity to the category query, and give the model only the best window or windows. The query uses only information available at inference time --- the category name or its CUAD question, never the answer --- so there is no leakage.

Before training anything, we measured the \emph{retrieval hit-rate}: how often does the selected window actually contain the gold clause? Figure~\ref{fig:ceiling} shows what the measurement is doing and why the answer is a hard bound. Whatever this number is, it is an upper bound on the downstream model's recall, because a clause sitting in a window the retriever never selects cannot be found by any model placed after it. The best strategy we found (category name as query, top-3 windows) reached only \textbf{64.0\%} (Table~\ref{tab:retrieval}). Interestingly, the CUAD question text did worse than the bare category name, because the questions share so much boilerplate that all the category queries end up nearly identical. A recall ceiling of 0.64 sits below the baseline's overall performance ($\approx$0.78), so a retrieval-windowed transformer was going to lose before it was ever trained. The measurement took seconds and saved us a wasted training run.

\begin{figure}[ht]
\centering
\begin{tikzpicture}[x=1mm, y=1mm, font=\scriptsize,
  wbox/.style={draw=gray!55, fill=gray!8, minimum width=15mm, minimum height=7mm,
               inner sep=0pt, rounded corners=1pt},
  sel/.style={draw=blue!65, fill=blue!20, minimum width=15mm, minimum height=7mm,
              inner sep=0pt, rounded corners=1pt, line width=0.7pt}
]
% ---- query --------------------------------------------------------------
\node[draw=gray!60, fill=orange!12, rounded corners=1pt, inner sep=2.5pt] (q) at (67,22)
  {category query $q_c$ = ``Governing Law''};
\draw[-{Stealth[length=2mm]}, gray!65] (67,17.6) -- (67,11.2);
\node[anchor=west, font=\tiny, text=gray!45!black] at (70,14.4)
  {TF--IDF cosine similarity, computed for every window};

% ---- windows ------------------------------------------------------------
\foreach \i/\x/\s in {1/0/0.12, 3/34/0.08, 5/68/0.19, 6/85/0.05, 8/119/0.10} {
  \node[wbox] (w\i) at (\x+7.5,7) {$w_{\i}$};
  \node[font=\tiny, text=gray!45!black] at (\x+7.5,1.5) {\s};
}
\foreach \i/\x/\s in {2/17/0.31, 4/51/0.27, 7/102/0.22} {
  \node[sel] (w\i) at (\x+7.5,7) {$w_{\i}$};
  \node[font=\tiny, text=blue!45!black] at (\x+7.5,1.5) {\textbf{\s}};
}
% gold-bearing window
\draw[riskhigh, line width=0.9pt, rounded corners=1pt] (67.2,2.8) rectangle (83.8,11.2);

% ---- legend + verdict ---------------------------------------------------
\node[anchor=west, font=\tiny] at (0,-4)
  {\textcolor{blue!65!black}{$\blacksquare$}~selected: top-$k{=}3$ windows sent to the model
   \qquad \textcolor{riskhigh}{$\square$}~window that actually holds the answer};
\node[anchor=west, align=left, font=\scriptsize] at (0,-10.5)
  {The gold window $w_5$ is \textbf{not} in the top~3, so this (contract, category) pair is
   \textbf{unreachable}:\\ no model downstream of the retriever can ever recover it. Repeated over
   all positive pairs, this yields a hit-rate of 64.0\%.};

% ---- ceiling bar --------------------------------------------------------
\fill[risklow!40] (0,-30) rectangle (89.6,-23);
\fill[gray!18]    (89.6,-30) rectangle (140,-23);
\draw[gray!60]    (0,-30) rectangle (140,-23);
\draw[gray!60]    (89.6,-30) -- (89.6,-23);
\node[font=\tiny] at (44.8,-26.5) {recall still attainable: \textbf{64.0\%}};
\node[font=\tiny, text=gray!45!black] at (114.8,-26.5) {lost before training starts: 36.0\%};
\draw[riskhigh, dashed, line width=0.8pt] (108.5,-32) -- (108.5,-20.5);
\node[anchor=south, font=\tiny, text=riskhigh!70!black] at (108.5,-20.2)
  {full-document TF--IDF baseline already scores 0.775};
\node[anchor=north, font=\tiny, text=gray!45!black] at (0,-30.6) {0\%};
\node[anchor=north, font=\tiny, text=gray!45!black] at (140,-30.6) {100\%};
\end{tikzpicture}
\caption{Why the retrieval ceiling settled the design. \emph{Top:} the query scores every window and only the top-$k$ go forward; whenever the gold-bearing window is not among them, the pair is lost no matter how good the downstream model is. \emph{Bottom:} the resulting 64.0\% ceiling sits \emph{below} what the cheap full-document baseline already achieves (0.775, dashed line), so retrieval could not have won --- which is why it was dropped before any training time was spent on it.}
\label{fig:ceiling}
\end{figure}

\begin{table}[ht]
\centering
\caption{Retrieval hit-rate on positive (contract, category) pairs --- the recall ceiling a retrieval-based design would impose.}
\label{tab:retrieval}
\begin{tabular}{lcc}
\toprule
\textbf{Query strategy} & \textbf{Top-$k$} & \textbf{Hit-rate}\\
\midrule
CUAD question text        & 1 & 38.6\%\\
Category name             & 1 & 44.6\%\\
\textbf{Category name}    & \textbf{3} & \textbf{64.0\%}\\
Name + question           & 3 & 61.5\%\\
\bottomrule
\end{tabular}
\end{table}

\subsection{Sliding-window max-pooling (multiple-instance learning)}
\label{sec:mil}
To get rid of the ceiling we trained a \emph{window-level} classifier instead: given the pair (category question, window), predict whether that window contains a clause of the category. Figure~\ref{fig:windows} shows how the windows are generated. A 2{,}000-character window advancing in 1{,}500-character steps leaves a 500-character overlap between neighbours, and that overlap is what guarantees no clause can fall between two windows and disappear: the median annotated clause is 196 characters (Section~\ref{sec:eda}), comfortably shorter than the overlap, so every clause lands wholly inside at least one window.

\begin{figure}[ht]
\centering
\begin{tikzpicture}[x=1mm, y=1mm, font=\scriptsize,
  dim/.style={{Stealth[length=1.6mm]}-{Stealth[length=1.6mm]}, gray!65!black, line width=0.5pt},
  win/.style={draw=blue!55, fill=blue!10, rounded corners=1pt},
  poswin/.style={draw=risklow!85, fill=risklow!20, rounded corners=1pt}]

% ---- overlap band (behind everything) -----------------------------------
\fill[orange!30, opacity=0.7] (15,-42) rectangle (20,0);

% ---- the document -------------------------------------------------------
\fill[gray!12] (0,0) rectangle (110,7);
\draw[gray!55] (0,0) rectangle (110,7);
\fill[riskhigh!60] (61,0) rectangle (63,7);
\node[anchor=west, text=gray!50!black, font=\tiny] at (2,3.5) {contract text};
\node[anchor=west] at (111.5,3.5) {$\cdots$};
\node[anchor=west, text=riskhigh!70!black, font=\tiny] at (70,11)
  {gold clause --- 196 chars, drawn to scale};
\draw[riskhigh!75, -{Stealth[length=1.6mm]}] (69.5,10.6) -- (63.4,7.4);

% ---- staircase of windows ----------------------------------------------
\foreach \i/\y in {0/-5, 1/-13, 2/-21} {
  \draw[win] (15*\i,\y-5) rectangle (15*\i+20,\y);
  \node[text=blue!45!black, font=\tiny] at (15*\i+10,\y-2.5) {$w_{\i}$};
}
\foreach \i/\y in {3/-29, 4/-37} {
  \draw[poswin] (15*\i,\y-5) rectangle (15*\i+20,\y);
  \node[text=risklow!35!black, font=\tiny] at (15*\i+10,\y-2.5) {$w_{\i}$};
  \fill[riskhigh!60] (61,\y-5) rectangle (63,\y);
}
\node[anchor=west, text=risklow!30!black, font=\tiny, align=left] at (86,-36)
  {both windows contain\\the clause in full};
\draw[risklow!70, -{Stealth[length=1.6mm]}] (85.5,-34.5) -- (64,-31.5);
\draw[risklow!70, -{Stealth[length=1.6mm]}] (85.5,-37.5) -- (64,-39.5);

% ---- dimensions ---------------------------------------------------------
\draw[dim] (0,-2.5) -- (20,-2.5);
\node[anchor=west, font=\tiny, text=gray!45!black] at (22,-2.5)
  {window $=$ 2{,}000 characters};
\draw[dim] (0,-10.5) -- (15,-10.5);
\node[anchor=west, font=\tiny, text=gray!45!black] at (22,-10.5)
  {stride $=$ 1{,}500 characters};
\draw[gray!55, dashed, line width=0.4pt] (15,-44) -- (15,7.5);
\draw[gray!55, dashed, line width=0.4pt] (20,-44) -- (20,7.5);
\node[anchor=north, align=center, font=\tiny, text=orange!50!black] at (17.5,-44)
  {overlap $=$ 500 chars, at every step};
\node[font=\normalsize, text=gray!55!black] at (67.5,-45.5) {$\vdots$};
\node[anchor=west, font=\tiny, text=gray!45!black] at (71,-45.5)
  {continuing to $w_{21}$ for a median-length contract};
\end{tikzpicture}
\caption{Sliding-window generation, drawn as a staircase so the geometry is explicit. Each 2{,}000-character window starts 1{,}500 characters after the previous one, so consecutive windows always share a 500-character overlap (orange column). Because that overlap is larger than the median annotated clause (196 characters), no clause can fall through the gap between two windows: here the clause lands inside \emph{both} $w_3$ and $w_4$ in full. This is what removes the recall ceiling of Figure~\ref{fig:ceiling}, and it is bought at the cost of scoring every window. A median-length contract produces 22 windows; the mean across the corpus is 35, and the longest contract yields 226.}
\label{fig:windows}
\end{figure}

Training labels come from the gold character offsets --- any window containing an annotated span is positive, and for each (contract, category) bag we sample two negative windows. Using gold offsets this way is ordinary supervision, not leakage: the labels shape the training targets, but they never touch the model's input, and nothing gold is used at inference. From the 408 training contracts this gives 53{,}662 window examples (39.2\% positive).

At inference the model scores \emph{every} window of the contract, and the document-level prediction is simply the maximum window probability --- the multiple-instance learning decision rule, drawn in Figure~\ref{fig:mil}. The document is a \emph{bag} of window instances; the bag is positive if any instance is positive; and $\max$ implements exactly that rule. Since the windows tile the whole document with overlap, every present clause is reachable, so the recall ceiling disappears by construction. The price is precision, and Figure~\ref{fig:mil} also shows why: the maximum is taken over every window, so a single spurious window anywhere in the document is enough to flip the whole contract to positive. With a median of 22 windows per contract and a mean of 35, the chance that some window crosses the 0.5 threshold by accident grows with document length --- long contracts are structurally more likely to produce a false positive than short ones, independently of what they actually contain.

\begin{figure}[ht]
\centering
\begin{tikzpicture}[x=1mm, y=1mm, font=\scriptsize,
  win/.style={draw=gray!55, fill=gray!10, rounded corners=1pt,
              minimum width=15mm, minimum height=6mm, inner sep=0pt},
  hot/.style={draw=riskhigh, fill=riskhigh!15, rounded corners=1pt,
              minimum width=15mm, minimum height=6mm, inner sep=0pt, line width=0.7pt},
  enc/.style={draw=blue!55, fill=blue!12, rounded corners=1pt,
              minimum width=24mm, minimum height=6mm, inner sep=0pt},
  arr/.style={-{Stealth[length=1.8mm]}, gray!60!black, line width=0.5pt}
]
% rows: w1, w2 (firing), w3, vdots, wn
\node[win] (w1) at (9,0)    {$w_1$};
\node[hot] (w2) at (9,-9)   {$w_2$};
\node[win] (w3) at (9,-18)  {$w_3$};
\node[font=\normalsize, text=gray!55!black] at (9,-25) {$\vdots$};
\node[win] (wn) at (9,-32)  {$w_n$};

\node[enc] (e1) at (45,0)   {DistilBERT};
\node[enc] (e2) at (45,-9)  {DistilBERT};
\node[enc] (e3) at (45,-18) {DistilBERT};
\node[font=\normalsize, text=gray!55!black] at (45,-25) {$\vdots$};
\node[enc] (en) at (45,-32) {DistilBERT};

\node (p1) at (77,0)   {$p_1 = 0.03$};
\node[text=riskhigh!70!black] (p2) at (77,-9) {$\mathbf{p_2 = 0.91}$};
\node (p3) at (77,-18) {$p_3 = 0.07$};
\node[font=\normalsize, text=gray!55!black] at (77,-25) {$\vdots$};
\node (pn) at (77,-32) {$p_n = 0.11$};

\foreach \a/\b in {w1/e1, w2/e2, w3/e3, wn/en} \draw[arr] (\a) -- (\b);
\foreach \a/\b in {e1/p1, e2/p2, e3/p3, en/pn} \draw[arr] (\a) -- (\b);

% max node
\node[draw=blue!65, fill=blue!15, rounded corners=1pt, minimum width=16mm,
      minimum height=9mm, inner sep=1pt, align=center] (mx) at (103,-16)
      {$\max\limits_i p_i$};
\foreach \p in {p1,p2,p3,pn} \draw[arr] (\p) -- (mx);

\node[anchor=west, align=left] (out) at (114,-16)
  {$p_{\text{doc}} = 0.91 \geq 0.5$\\[1pt] $\Rightarrow$ category \textbf{present}};
\draw[arr] (mx) -- (out);

% annotations
\node[anchor=north, align=center, font=\tiny, text=gray!45!black] at (60,-37)
  {one model, shared weights --- the input is the pair (category question $q_c$, window $w_i$)};
\end{tikzpicture}
\caption{The multiple-instance learning decision rule. Every window of the contract is scored independently against the same category question by one shared DistilBERT, and the document-level probability is the maximum over windows. Here only the shaded window $w_2$ actually contains the clause, and its score alone determines the document-level answer. One firing window is enough to call the category present --- which is what makes recall complete, and equally what makes precision fragile, since a single spurious window anywhere in a long contract produces a document-level false positive.}
\label{fig:mil}
\end{figure}

\subsection{The conjunctive ensemble}
When we compared the two models category by category, they turned out to be \emph{complementary}: the transformer wins on semantically subtle categories (for example \emph{No-Solicit of Employees} and \emph{Liquidated Damages}), while the baseline wins on keyword-heavy ones. We combine them with \textbf{parameter-free} rules --- AND ($\min$ of the two probabilities), OR ($\max$), and AVG (mean) --- chosen deliberately because they contain nothing tunable, so nothing can overfit to the test set. Figure~\ref{fig:ensemble_rule} lays out the decision. The AND rule, which reports a clause only when both models agree, filters out each model's uncorrelated false positives: for a spurious detection to survive, the lexical model and the windowed transformer must both make the same mistake on the same contract and category, and their mistakes are largely independent. That recovers exactly the precision that max-pooling loses.

\begin{figure}[ht]
\centering
\begin{tikzpicture}[x=1mm, y=1mm, font=\scriptsize,
  mod/.style={draw=gray!60, fill=gray!10, rounded corners=1pt, text width=30mm,
              align=center, minimum height=11mm, inner sep=1.5pt},
  rule/.style={draw=gray!60, fill=gray!8, rounded corners=1pt, text width=24mm,
               align=center, minimum height=8mm, inner sep=1.5pt},
  pick/.style={rule, draw=blue!70, fill=blue!18, line width=0.8pt},
  arr/.style={-{Stealth[length=1.8mm]}, gray!60!black, line width=0.5pt}
]
\node[draw=gray!60, fill=orange!12, rounded corners=1pt, text width=16mm, align=center,
      minimum height=9mm, inner sep=1.5pt] (doc) at (9,0) {contract\\text};

\node[mod] (bl) at (46,15) {TF--IDF $+$ LR\\{\tiny reads the \emph{whole} document}};
\node[mod] (tr) at (46,-15) {DistilBERT MIL\\{\tiny max-pools over windows}};
\draw[arr] (doc) -- (bl);
\draw[arr] (doc) -- (tr);

\node (pa) at (76,15)  {$p_{\text{tfidf}}$};
\node (pb) at (76,-15) {$p_{\text{trans}}$};
\draw[arr] (bl) -- (pa);
\draw[arr] (tr) -- (pb);

\coordinate (j) at (86,0);
\draw[gray!60!black, line width=0.5pt] (pa) -- (86,15) -- (j);
\draw[gray!60!black, line width=0.5pt] (pb) -- (86,-15) -- (j);
\fill[gray!60!black] (j) circle (0.6mm);

\node[pick] (rand) at (110,15)  {\textbf{AND}\\[1pt]$\min(p_a,p_b) \geq 0.5$};
\node[rule] (ravg) at (110,0)   {AVG\\[1pt]$\tfrac{1}{2}(p_a{+}p_b) \geq 0.5$};
\node[rule] (ror)  at (110,-15) {OR\\[1pt]$\max(p_a,p_b) \geq 0.5$};
\foreach \r in {rand,ravg,ror} \draw[arr] (j) -- (\r);

\node[anchor=west, font=\tiny, text=blue!45!black] at (124,15) {\textbf{0.776}};
\node[anchor=west, font=\tiny, text=gray!45!black]  at (124,0)  {0.718};
\node[anchor=west, font=\tiny, text=gray!45!black]  at (124,-15){0.696};
\node[anchor=south, font=\tiny, text=gray!45!black, align=center] at (130,20)
  {held-out\\micro-F1};

\node[anchor=north, align=center, font=\tiny, text=gray!45!black] at (67,-24)
  {All three rules are \textbf{parameter-free}: there is nothing to fit, so the test set cannot tune them.
   AND fires only when \emph{both} models\\ cross the threshold, so a false positive has to be made
   twice, independently, to survive.};
\end{tikzpicture}
\caption{The conjunctive ensemble. The two presence models see the contract in fundamentally different ways --- lexical evidence over the entire document versus semantic evidence over overlapping windows --- and their errors are therefore largely uncorrelated. Requiring agreement (AND) discards the false positives that max-pooling introduces. The micro-F1 values on the right are the held-out results reported in Chapter~6, shown here only to indicate which rule the analysis selects.}
\label{fig:ensemble_rule}
\end{figure}

\section{Stage 2B: Span Extraction}
Here the transformer has no substitute --- a bag-of-words model cannot point at character positions. We fine-tune \texttt{DistilBertForQuestionAnswering}: the input is (CUAD question, window text) and the model predicts start and end token positions. Character offsets from \texttt{CUAD\_v1.json} are converted to token labels through the tokenizer's offset mapping.

Two different window sizes are in play in this thesis, and it is worth being explicit about the difference, since mixing them up makes the whole design look inconsistent. The 2{,}000-character window of Section~\ref{sec:mil} is a \emph{presence} window: it tiles the document blindly, because at that point we do not know where anything is. The window used here is a \emph{focus} window: 1{,}200 characters, and it is not tiled at all but re-centred on the gold answer, positioned so the answer begins roughly 150 characters in. Figure~\ref{fig:focus} draws the two on the same scale. The reason for the smaller, re-centred window is purely mechanical: it guarantees the target span survives truncation at the model's 320-token limit, which a blind 2{,}000-character window does not --- an answer sitting at the tail of a presence window would simply be cut off, and the model would be trained on an example whose answer is not present. Note also that this stage has no length ceiling at all: extraction happens \emph{inside} a window that Stage 2A has already flagged as relevant.

\begin{figure}[ht]
\centering
\begin{tikzpicture}[x=1mm, y=1mm, font=\scriptsize,
  dim/.style={{Stealth[length=1.6mm]}-{Stealth[length=1.6mm]}, gray!65!black, line width=0.5pt}]

% ---- presence window (Stage 2A) ----------------------------------------
\fill[blue!10] (0,0) rectangle (120,8);
\draw[blue!55] (0,0) rectangle (120,8);
\fill[riskhigh!35] (54,0) rectangle (65.8,8);
\draw[riskhigh]    (54,0) rectangle (65.8,8);
\node[anchor=west, align=left, font=\tiny, text=gray!45!black] at (123,4)
  {\textbf{Stage 2A}\\presence window\\2{,}000 chars, tiled};
\node[anchor=south, font=\tiny, text=riskhigh!70!black] at (59.9,11.6) {gold answer};
\draw[riskhigh!75, -{Stealth[length=1.6mm]}] (59.9,11.2) -- (59.9,8.6);

% ---- focus window (Stage 2B) -------------------------------------------
\fill[risklow!25] (45,-16) rectangle (117,-8);
\draw[risklow!85] (45,-16) rectangle (117,-8);
\fill[riskhigh!35] (54,-16) rectangle (65.8,-8);
\draw[riskhigh]    (54,-16) rectangle (65.8,-8);
\node[anchor=west, align=left, font=\tiny, text=gray!45!black] at (123,-12)
  {\textbf{Stage 2B}\\focus window\\1{,}200 chars, re-centred};

% ---- alignment guides + offset -----------------------------------------
\draw[gray!55, dashed, thin] (45,-16) -- (45,-2);
\draw[gray!55, dashed, thin] (54,-16) -- (54,-1);
\draw[dim] (45,-4.2) -- (54,-4.2);
\node[anchor=west, font=\tiny, text=gray!45!black] at (56,-4.2)
  {150 chars --- the answer never sits on the boundary};

% ---- footer ------------------------------------------------------------
\node[anchor=north, align=center, font=\tiny, text=gray!45!black] at (60,-21)
  {$\text{focus start} = \max(0,\ \text{answer start} - 150)$\quad ---\quad
   1{,}200 chars fits inside the 320-token limit; 2{,}000 does not.};
\end{tikzpicture}
\caption{The two window sizes, drawn to the same scale. Stage 2A tiles the document with 2{,}000-character \emph{presence} windows to find out \emph{whether} a clause is there. Stage 2B then trains on a narrower 1{,}200-character \emph{focus} window re-centred on the answer, so the span is guaranteed to survive tokenizer truncation. The two are not alternatives and never compete: the first locates a relevant region, the second reads inside it.}
\label{fig:focus}
\end{figure}

\section{Stage 1: Plain-Language Summarization}
CUAD gives us clauses but no plain-English rewrites to learn from, so we wrote our own seed: \textbf{41 gold templates}, one carefully phrased single-sentence summary per category, framed from the weaker party's perspective. For \emph{Anti-Assignment}, for instance, the template reads ``You cannot transfer this contract to someone else without permission.'' Every extracted training clause is paired with its category's template, which yields 11{,}156 training pairs, and FLAN-T5-small is fine-tuned on the prompt ``\texttt{summarize in plain English: <clause>}''.

We evaluate with ROUGE-L \emph{and} with a side-by-side zero-shot-versus-fine-tuned comparison on the same clauses. The double evaluation is deliberate. With only 41 unique target sentences, a model can earn a high ROUGE score just by learning to classify the clause and emit its template word for word, without doing any real summarization. Chapter 6 shows that this is exactly what happens, and we report it as such.

\section{Risk Taxonomy}
Each of the 41 categories carries a fixed risk level --- \textcolor{riskhigh}{\textbf{High}}, \textcolor{riskmed}{\textbf{Medium}}, or \textcolor{risklow}{\textbf{Low}} --- scored from the weaker party's perspective, together with a one-line reason. The taxonomy (8 High, 22 Medium, 11 Low; Figure~\ref{fig:riskdist}) is a transparent lookup rather than a learned score: it can be audited line by line, its reasoning is written down per category, and it cannot drift silently. The High categories fall into three buckets: money you could owe without limit (\emph{Uncapped Liability}, \emph{Liquidated Damages}, \emph{Minimum Commitment}, \emph{Most Favored Nation}); losing something you cannot get back (\emph{IP Ownership Assignment}, \emph{Irrevocable/Perpetual License}); and losing the freedom to earn later (\emph{Non-Compete}, \emph{Exclusivity}). The full clause-by-clause table with reasons is in Appendix~\ref{app:taxonomy}.

\begin{figure}[ht]
\centering
\includegraphics[width=0.62\textwidth]{risk_distribution.png}
\caption{Distribution of the 41 CUAD clause categories across the three risk levels.}
\label{fig:riskdist}
\end{figure}

\section{Leakage Discipline}
With a dataset this small, leakage is the most dangerous way to fool ourselves, so we kept one explicit rule the whole way through: \emph{labels may shape training targets, but they may never shape the input a model sees or any inference-time selection}. In practice this meant: the split is contract-level; retrieval queries use only the category or question; gold offsets label training windows but never build inputs or pick inference windows; and the ensemble rules are parameter-free, so the test set cannot tune them.

One consequence is visible in Figure~\ref{fig:focus} and deserves stating outright, because it is the one place where the rule looks like it might be broken. The focus window is positioned using the gold answer offset, which is a label. That is legitimate only because it happens exclusively at \emph{training} time: it decides which text becomes a training example, never what the model reads at inference. At inference there is no gold offset, and Stage 2B simply runs over the presence window that Stage 2A selected. If re-centring were applied at inference the evaluation would be meaningless, since the model would be handed a window built around the answer it is supposed to find.

\section{End-to-End Data Flow}
Figure~\ref{fig:pipeline} names the stages; Figure~\ref{fig:dataflow} adds what actually moves between them. The shapes shown are for one contract of median length passing through the deployed system.

\begin{figure}[ht]
\centering
\begin{tikzpicture}[x=1mm, y=1mm, font=\scriptsize,
  st/.style={draw=blue!55, fill=blue!8, rounded corners=2pt, text width=40mm,
             align=center, minimum height=8mm, inner sep=1.5pt},
  ad/.style={st, draw=orange!75, fill=orange!15},
  arr/.style={-{Stealth[length=2mm]}, gray!60!black, line width=0.6pt},
  op/.style={font=\tiny, text=gray!45!black, anchor=west},
  dat/.style={font=\tiny, anchor=west, text width=76mm}
]
\node[st] (s1) at (24,0)   {Contract text};
\node[st] (s2) at (24,-16) {Sliding windows};
\node[st] (s3) at (24,-32) {Presence scoring};
\node[st] (s4) at (24,-48) {Presence vector};
\node[st] (s5) at (24,-64) {Extracted spans};
\node[ad] (s6) at (24,-80) {Risk $+$ plain English};
\node[ad] (s7) at (24,-96) {Risk-prioritized report};

\foreach \a/\b in {s1/s2, s2/s3, s3/s4, s4/s5, s5/s6, s6/s7} \draw[arr] (\a) -- (\b);

\node[op] at (26,-8)  {windowing: 2{,}000 chars, stride 1{,}500};
\node[op] at (26,-24) {$\times$ 41 category questions; DistilBERT $+$ TF--IDF};
\node[op] at (26,-40) {$\max$ over windows, then AND with TF--IDF};
\node[op] at (26,-56) {DistilBERT-QA, flagged categories only};
\node[op] at (26,-72) {taxonomy lookup $+$ FLAN-T5-small};
\node[op] at (26,-88) {sort by risk level, then confidence};

\node[dat] at (50,0)   {\texttt{str} --- 33{,}143 characters (median)};
\node[dat] at (50,-16) {\texttt{list[str]} --- $n = 22$ windows (median; mean 35, max 226)};
\node[dat] at (50,-32) {\texttt{float[n,\,41]} window scores $+$ \texttt{float[41]} document scores};
\node[dat] at (50,-48) {\texttt{bool[41]} --- about 13 categories flagged present};
\node[dat] at (50,-64) {\texttt{list[(category,\,char\_start,\,char\_end,\,text)]}};
\node[dat] at (50,-80) {\texttt{list[(category,\,High\,$|$\,Med\,$|$\,Low,\,reason,\,sentence)]}};
\node[dat] at (50,-96) {grouped by risk level, rendered in \emph{Clausify} (Chapter~7)};

\draw[gray!45, dashed, line width=0.4pt] (47,4) -- (47,-100);
\node[anchor=south, font=\tiny, text=gray!45!black] at (24,5) {\textbf{stage}};
\node[anchor=south west, font=\tiny, text=gray!45!black] at (50,5) {\textbf{data leaving that stage}};
\end{tikzpicture}
\caption{The pipeline of Figure~\ref{fig:pipeline} annotated with the object that crosses each stage boundary, for one median-length contract. The narrowing from \texttt{float[n,\,41]} window scores to a \texttt{bool[41]} presence vector is where max-pooling and the AND rule act; everything downstream operates only on the roughly 13 categories that survive, which is what keeps inference fast enough for an interactive application.}
\label{fig:dataflow}
\end{figure}
